\paragraph{Interference and Final Measurement}\label{ex:interference-and-final-measurement}

After the oracle has been applied, the computation qubit remains in the state \ket{-}, while the data qubit contains the information about the function encoded as a pattern of relative phases.

\highspace
The state of the two qubits is:
\begin{equation*}
    \ket{\psi_2} = \frac{1}{\sqrt{2}} \left( (-1)^{f(0)}\ket{0} + (-1)^{f(1)}\ket{1} \right) \otimes \ket{-}
\end{equation*}
\textcolor{Red2}{\faIcon{exclamation-triangle}} At this point, the computation is almost complete. However, there is still a \textbf{problem}: \hl{relative phases cannot be observed directly by measurement.} If we were to measure the data qubit immediately, we would obtain:
\begin{equation*}
    P(0) = P(1) = \frac{1}{2}
\end{equation*}
Regardless of the values of $f(0)$ and $f(1)$. The information is present in the quantum state, but it is hidden in the interference between the two amplitudes rather than in their magnitudes. Therefore, to reveal this information, \hl{Deutsch's algorithm applies one final Hadamard gate to the data qubit.}

\highspace
\begin{flushleft}
    \textcolor{Green3}{\faIcon{check-circle} \textbf{Solution: Applying the Final Hadamard Gate}}
\end{flushleft}
The final operation of the algorithm is:
\begin{equation*}
    H \otimes I
\end{equation*}
Which \textbf{acts only on the data qubit}. Recall the Hadamard transformations:
\begin{equation*}
    H\ket{+} = \ket{0}, \quad H\ket{-} = \ket{1}
\end{equation*}
Therefore, if we can determine whether the data qubit is in the state \ket{+} or \ket{-}, the Hadamard immediately converts this information into one of the computational basis states.
\begin{equation}
    \begin{aligned}
        \ket{\psi_3} &= (H \otimes I)\ket{\psi_2} \\
        &= \frac{1}{\sqrt{2}} \left( (-1)^{f(0)}H\ket{0} + (-1)^{f(1)}H\ket{1} \right) \otimes \ket{-} \\
        &= \frac{1}{2} \left[ (-1)^{f(0)} \left(\ket{0} + \ket{1}\right) + (-1)^{f(1)} \left(\ket{0} - \ket{1}\right) \right] \otimes \ket{-} \\
        &= \frac{1}{2} \left[ \left((-1)^{f(0)} + (-1)^{f(1)}\right)\ket{0} + \left((-1)^{f(0)} - (-1)^{f(1)}\right)\ket{1} \right] \otimes \ket{-}
    \end{aligned}
\end{equation}
This is the \textbf{essence of quantum interference}: \hl{the Hadamard combines the amplitudes of the two branches so that they either reinforce or cancel each other.} From the previous section, we know that the data qubit is in the state \ket{+} if $f(0) = f(1)$ and in the state \ket{-} if $f(0) \neq f(1)$. Therefore, after applying the final Hadamard gate, we have:
\begin{itemize}
    \item \important{Constant Functions}. Every constant function leaves the data qubit in the state $\pm \ket{+}$. Applying the Hadamard gives:
    \begin{equation*}
        H\left(\pm \ket{+}\right) = \pm \ket{0}
    \end{equation*}
    The minus sign is merely a global phase and therefore has no physical effect. Consequently, \hl{if the function is constant, the final measurement of the data qubit will always yield $\ket{0}$.} The probability of measuring $\ket{0}$ is 100\%, while the probability of measuring $\ket{1}$ is 0\% (because the amplitude of $\ket{1}$ is zero). This is a \textbf{deterministic outcome} that allows us to conclude that the function is constant with certainty.


    \item \important{Balanced Functions}. Similarly, every balanced function leaves the data qubit in the state $\pm \ket{-}$. Applying the Hadamard gives:
    \begin{equation*}
        H\left(\pm \ket{-}\right) = \pm \ket{1}
    \end{equation*}
    Again, the minus sign is a global phase and has no physical effect. Therefore, \hl{if the function is balanced, the final measurement of the data qubit will always yield $\ket{1}$.} The probability of measuring $\ket{1}$ is 100\%, while the probability of measuring $\ket{0}$ is 0\% (because the amplitude of $\ket{0}$ is zero). This is also a \textbf{deterministic outcome} that allows us to conclude that the function is balanced with certainty.
\end{itemize}
\textcolor{Green3}{\faIcon{question-circle} \textbf{Understanding the Interference.}} To understand how the Hadamard gate causes interference, recall the definition of interference on \autopageref{def:interference}. Although it is placed in the context of quantum annealing, it is not necessary to understand the details of that algorithm to understand interference. 
What is important are the two images in the example, which show a similar effect created by the Hadamard gate.

\highspace
The Hadamard gate effectively adds and subtracts the amplitudes of the two computational basis states. Recall that:
\begin{equation*}
    H = \frac{1}{\sqrt{2}}\begin{pmatrix}
        1 & 1 \\
        1 & -1
    \end{pmatrix}
\end{equation*}
If the amplitudes have the \textbf{same sign}, as in:
\begin{equation*}
    \ket{+} = \frac{1}{\sqrt{2}}\left(\ket{0} + \ket{1}\right)
\end{equation*}
They interfere constructively in the first output component and destructively in the second:
\begin{equation*}
    H\ket{+} = \frac{1}{\sqrt{2}}\begin{pmatrix}
        1 & 1 \\
        1 & -1
    \end{pmatrix}\frac{1}{\sqrt{2}}\begin{pmatrix}
        1 \\
        1
    \end{pmatrix} = \frac{1}{2}\begin{pmatrix}
        2 \\
        0
    \end{pmatrix} = \ket{0}
\end{equation*}
Conversely, if the amplitudes have opposite signs, as in:
\begin{equation*}
    \ket{-} = \frac{1}{\sqrt{2}}\left(\ket{0} - \ket{1}\right)
\end{equation*}
The interference occurs in the opposite way:
\begin{equation*}
    H\ket{-} = \frac{1}{\sqrt{2}}\begin{pmatrix}
        1 & 1 \\
        1 & -1
    \end{pmatrix}\frac{1}{\sqrt{2}}\begin{pmatrix}
        1 \\
        -1
    \end{pmatrix} = \frac{1}{2}\begin{pmatrix}
        0 \\
        2
    \end{pmatrix} = \ket{1}
\end{equation*}
Thus, the \hl{Hadamard converts the \textbf{relative phase} introduced by the oracle into a measurable computational-basis outcome.} This is why quantum interference is essential: without it, the phase information would remain inaccessible.

\highspace
\begin{flushleft}
    \textcolor{Green3}{\faIcon{check-circle} \textbf{Final Decision Rule}}
\end{flushleft}
The algorithm now performs a \textbf{measurement} of the \textbf{data qubit}. The result is deterministic:
\begin{itemize}
    \item Measuring $\ket{0}$ indicates that the function is constant.
    \item Measuring $\ket{1}$ indicates that the function is balanced.
\end{itemize}
Notice that the \textbf{computation qubit} is \textbf{never measured}. Its sole purpose was to enable phase kickback inside the oracle. After transferring the function values into the relative phases of the data qubit, it no longer carries any useful information for the algorithm.

\highspace
In summary, the complete evolution of \definition{Deutsch's Algorithm} can now be summarized as:
\begin{equation}
    \begin{aligned}
        \ket{0,1} &\xrightarrow{H \otimes H} \ket{+,-} \\
        &\xrightarrow{U_f} \frac{1}{\sqrt{2}} \left( (-1)^{f(0)}\ket{0} + (-1)^{f(1)}\ket{1} \right) \otimes \ket{-} \\
        &\xrightarrow{H \otimes I} \begin{cases}
            \ket{0, -} & \text{if } f \text{ is constant} \\
            \ket{1, -} & \text{if } f \text{ is balanced}
        \end{cases}
    \end{aligned}
\end{equation}
And all the states:
\begin{equation*}
    \begin{aligned}
        \ket{\psi_{0}} &= \ket{0,1} \\
        \ket{\psi_{1}} &= \ket{+,-} \\
        \ket{\psi_{2}} &= \frac{1}{\sqrt{2}} \left( (-1)^{f(0)}\ket{0} + (-1)^{f(1)}\ket{1} \right) \otimes \ket{-} \\
        \ket{\psi_{3}} &= \begin{cases}
            \ket{0, -} & \text{if } f \text{ is constant} \\
            \ket{1, -} & \text{if } f \text{ is balanced}
        \end{cases} \quad \text{(up to a global phase)}
    \end{aligned}
\end{equation*}
\begin{figure}[!htp]
    \centering
    \begin{quantikz}
        \lstick{\ket{0}}\slice{$\ket{\psi_{0}}$}    & \gate{H}\slice{$\ket{\psi_{1}}$}  & \ctrl{1}\slice{$\ket{\psi_{2}}$}  & \gate{H}\slice{$\ket{\psi_{3}}$}  & \meter{}  & \\
        \lstick{\ket{1}}                            & \gate{H}                          & \gate{U_f}                        &                                   &           &
    \end{quantikz}
    \caption{Deutsch's Algorithm. The final measurement of the data qubit reveals whether the function is constant or balanced.}
\end{figure}

\newpage

\begin{examplebox}[: $f(x) = 0$]
    Let's say we have a constant function $f(x) = 0$:
    \begin{equation*}
        f(0) = f(1) = 0
    \end{equation*}
    Therefore, the state after the oracle is:
    \begin{align*}
        \ket{\psi_2} &= \frac{1}{\sqrt{2}} \left( (-1)^{f(0)}\ket{0} + (-1)^{f(1)}\ket{1} \right) \otimes \ket{-} \\
        &= \frac{1}{\sqrt{2}} \left( (-1)^{0}\ket{0} + (-1)^{0}\ket{1} \right) \otimes \ket{-} \\
        &= \frac{1}{\sqrt{2}} \left( \ket{0} + \ket{1} \right) \otimes \ket{-} \\
        &= \ket{+,-}
    \end{align*}
    Applying the final Hadamard gate gives:
    \begin{equation*}
        \ket{\psi_3} = (H \otimes I)\ket{\psi_2} = (H \otimes I)\ket{+,-} = H\ket{+} \otimes \ket{-} = \ket{0,-}
    \end{equation*}
    Thus the measured data qubit is $\ket{0}$, indicating that the function is \textbf{constant}.
\end{examplebox}

\highspace
\begin{examplebox}[: $f(x) = 1$]
    Let's say we have a constant function $f(x) = 1$:
    \begin{equation*}
        f(0) = f(1) = 1
    \end{equation*}
    Therefore, the state after the oracle is:
    \begin{align*}
        \ket{\psi_2} &= \frac{1}{\sqrt{2}} \left( (-1)^{f(0)}\ket{0} + (-1)^{f(1)}\ket{1} \right) \otimes \ket{-} \\
        &= \frac{1}{\sqrt{2}} \left( (-1)^{1}\ket{0} + (-1)^{1}\ket{1} \right) \otimes \ket{-} \\
        &= -\frac{1}{\sqrt{2}} \left( \ket{0} + \ket{1} \right) \otimes \ket{-} \\
        &= -\ket{+,-}
    \end{align*}
    Applying the final Hadamard gate gives:
    \begin{equation*}
        \ket{\psi_3} = (H \otimes I)\ket{\psi_2} = (H \otimes I)(-\ket{+,-}) = -H\ket{+} \otimes \ket{-} = -\ket{0,-}
    \end{equation*}
    Thus the measured data qubit is $\ket{0}$, indicating that the function is also \textbf{constant}. The minus sign is a global phase and has no physical effect on the measurement outcome.
\end{examplebox}

\highspace
\begin{examplebox}[: $f(x) = x$]
    Let's say we have a balanced function $f(x) = x$:
    \begin{equation*}
        f(0) = 0, \quad f(1) = 1
    \end{equation*}
    Therefore, the state after the oracle is:
    \begin{align*}
        \ket{\psi_2} &= \frac{1}{\sqrt{2}} \left( (-1)^{f(0)}\ket{0} + (-1)^{f(1)}\ket{1} \right) \otimes \ket{-} \\
        &= \frac{1}{\sqrt{2}} \left( (-1)^{0}\ket{0} + (-1)^{1}\ket{1} \right) \otimes \ket{-} \\
        &= \frac{1}{\sqrt{2}} \left( \ket{0} - \ket{1} \right) \otimes \ket{-} \\
        &= \ket{-,-}
    \end{align*}
    Applying the final Hadamard gate gives:
    \begin{equation*}
        \ket{\psi_3} = (H \otimes I)\ket{\psi_2} = (H \otimes I)\ket{-,-} = H\ket{-} \otimes \ket{-} = \ket{1,-}
    \end{equation*}
    Thus the measured data qubit is $\ket{1}$, indicating that the function is \textbf{balanced}.
\end{examplebox}

\highspace
\begin{examplebox}[: $f(x) = \overline{x}$]
    Let's say we have a balanced function $f(x) = \overline{x}$:
    \begin{equation*}
        f(0) = 1, \quad f(1) = 0
    \end{equation*}
    Therefore, the state after the oracle is:
    \begin{align*}
        \ket{\psi_2} &= \frac{1}{\sqrt{2}} \left( (-1)^{f(0)}\ket{0} + (-1)^{f(1)}\ket{1} \right) \otimes \ket{-} \\
        &= \frac{1}{\sqrt{2}} \left( (-1)^{1}\ket{0} + (-1)^{0}\ket{1} \right) \otimes \ket{-} \\
        &= -\frac{1}{\sqrt{2}} \left( -\ket{0} + \ket{1} \right) \otimes \ket{-} \\
        &= -\ket{-,-}
    \end{align*}
    Applying the final Hadamard gate gives:
    \begin{equation*}
        \ket{\psi_3} = (H \otimes I)\ket{\psi_2} = (H \otimes I)(-\ket{-,-}) = -H\ket{-} \otimes \ket{-} = -\ket{1,-}
    \end{equation*}
    Thus the measured data qubit is $\ket{1}$, indicating that the function is also \textbf{balanced}. The minus sign is a global phase and has no physical effect on the measurement outcome.
\end{examplebox}